\section{Methods}\label{sec:methods}

\subsection{A latent Gaussian model for event occurrence}\label{sec:methods-model}

Let $\mathcal{S}$ be the set of $n$ grid cells, $c(s)$ the country and $b(s)$
the biome of cell $s$, and $t = 1, \dots, T$ the weeks. Conditionally on a set
of latent effects, the outcomes are independent Bernoulli variables,
$Y_{st} \mid \pi_{st} \sim \mathrm{Bernoulli}(\pi_{st})$, with
\begin{equation}\label{eq:model}
\logit \pi_{st}
= \bx_{st}^\top \bbeta
+ \delta\, h_{st}
+ \sum_{k \in \mathcal{K}} z_{k,st}\, \beta_k(s)
+ \bz_{c(s),t}^\top \bm{\alpha}
+ u_{c(s)} + w(s) + e_s + v(s,t).
\end{equation}
Here $h_{st}$ is the recent event history of Section~\ref{sec:data-history},
so $\delta$ measures state dependence (Section~\ref{sec:methods-persistence}).
The vector $\bx_{st}$ holds the cell and cell-week covariates with fixed effects.
In the application these are the other outcome's event history, population,
malaria risk, the IPC phase 3+ population, river length and lake area, the climate
anomalies and climatologies of Section~\ref{sec:data-cov}, biome effects and year
effects. The variables $z_{k,st}$, $k \in \mathcal{K}$, are climate anomalies whose
effects $\beta_k(s)$ vary in space (Section~\ref{sec:methods-svc}); in the
application $\mathcal{K}$ contains maximum temperature only. The vector
$\bz_{ct}$ holds the national indicators with effects $\bm{\alpha}$
(Section~\ref{sec:methods-confounding}). GDP per capita is one of them: it is
national GDP divided by national population and enters as its logarithm. The
country effects $u_c \sim \Normal(0, \sigma_u^2)$ are independent, $w$ is a
Mat\'ern field on the sphere (Section~\ref{sec:methods-spde}), the cell effects
$e_s \sim \Normal(0, \sigma_e^2)$ are independent, and $v$ is a space-time field
(Section~\ref{sec:methods-st}). The four latent terms capture persistent
\emph{recorded} risk at four scales. Each mixes occurrence with reporting
(Section~\ref{sec:data-acled}): national institutions and reporting regimes ($u$),
regional hotspots that cross borders ($w$), single cells such as a town, a mine or
a border crossing ($e$), and slowly evolving episodes ($v$). The cell effect is
needed because a field with $\nu = 1$ and a range of several cells cannot represent
risk that differs between adjacent cells. Without it, such risk would be absorbed
by $w$ or by the history term $\delta h_{st}$ (Section~\ref{sec:methods-persistence}).

Missing covariate values are kept as missing through the construction of the
features. In particular, a cell-week without an IPC assessment has a missing IPC
value, not a food-insecure population of zero. After standardisation, missing values
are set to the training-period mean, and a missingness indicator is added as a
further fixed effect for each covariate group: one indicator for each national
indicator and one for IPC. Without these indicators, the imputed value would be
confounded with the fact of not being assessed, which is itself associated with
conflict. The boosting models receive the missing values as such, since XGBoost
handles them natively.

Model~\eqref{eq:model} is a latent Gaussian model \citep{rue2009inla}. Stacking all
latent quantities into
$\bm{\xi} = (\bbeta, \delta, \bgamma, \bm{\alpha}, \bu, \bw, \bm{e}, \bmeta, \bv)$, the
linear predictor is $\bm{\eta} = \bm{o} + \bA \bm{\xi}$ for a sparse design
matrix $\bA$ and an optional offset $\bm{o}$, and
$\bm{\xi} \mid \btheta \sim \Normal\{\bm{0}, \bQ(\btheta)^{-1}\}$ with a sparse,
block-diagonal precision whose blocks are listed in Table~S1 of the Supplementary Material
\citep{rue2005gmrf}. Fixed effects receive vague $\Normal(0, 10^4)$ priors and are
integrated out together with the random effects.

\subsection{Mat\'ern fields on the sphere via a finite-volume SPDE}\label{sec:methods-spde}

Following \citet{lindgren2011spde}, a Gaussian field $w$ on $\sphere$ with
Mat\'ern-like covariance is defined as the stationary solution of
\begin{equation}\label{eq:spde}
\tau\,(\kappa^2 - \Delta_{\sphere})\, w(s) = \mathcal{W}(s), \qquad s \in \sphere,
\end{equation}
where $\Delta_{\sphere}$ is the Laplace--Beltrami operator and $\mathcal{W}$ is
Gaussian white noise \citep[see also][]{lindgren2022spde}. This is the $\alpha = 2$
case, which corresponds to Mat\'ern smoothness $\nu = 1$ in the planar limit. The
field then has practical range $\rho = \sqrt{8}/\kappa$ and marginal variance
$\sigma^2 = 1/(4\pi\kappa^2\tau^2)$, provided $\rho$ is small relative to the
Earth's radius $R$. We parameterise the field by $(\rho, \sigma)$, in kilometres and
logit units.

The grid cells are natural control volumes, so instead of a triangulated mesh
we discretise \eqref{eq:spde} by finite volumes on the longitude--latitude
lattice itself \citep[for finite-volume SPDE discretisations on regular grids,
see][]{fuglstad2015exploring}. Integrating \eqref{eq:spde} over cell $A_i$ and
applying the divergence theorem with a two-point flux approximation gives
\[
\tau\Bigl\{\kappa^2 a_i w_i + \sum_{j \sim i} \frac{\ell_{ij}}{d_{ij}}(w_i - w_j)\Bigr\}
= \mathcal{W}(A_i) \sim \Normal(0, a_i),
\]
where $a_i$ is the area of cell $i$, $j \sim i$ ranges over the cells sharing an
edge with $i$, $\ell_{ij}$ is the length of the shared edge and $d_{ij}$ the
great-circle distance between the centres. For a grid with spacings
$\Delta\lambda, \Delta\varphi$ and centre latitude $\varphi_i$,
\begin{align*}
a_i &= R^2 \Delta\lambda\,\{\sin(\varphi_i + \Delta\varphi/2) - \sin(\varphi_i - \Delta\varphi/2)\},\\
\ell_{ij}/d_{ij} &= \Delta\varphi / (\cos\varphi_i\, \Delta\lambda) \quad \text{(east--west neighbours)},\\
\ell_{ij}/d_{ij} &= \cos\varphi_{ij}\,\Delta\lambda / \Delta\varphi \quad \text{(north--south neighbours)},
\end{align*}
with $\varphi_{ij}$ the latitude of the shared edge. Writing
$\bC = \mathrm{diag}(a_i)$ and $\bG$ for the weighted graph Laplacian with edge
weights $\ell_{ij}/d_{ij}$, we get $\tau \bK \bw = \bC^{1/2}\bm{\epsilon}$ with
$\bK = \kappa^2 \bC + \bG$ and $\bm{\epsilon} \sim \Normal(\bm{0}, \bm{I})$, so
\begin{equation}\label{eq:Qs}
\bw \sim \Normal\{\bm{0}, \bQ_s^{-1}\}, \qquad
\bQ_s(\rho, \sigma) = \tau^2 \bK \bC^{-1} \bK,
\qquad
\log|\bQ_s| = 2n\log\tau + 2\log|\bK| - \textstyle\sum_i \log a_i .
\end{equation}
$\bQ_s$ has at most 13 non-zeros per row. The construction is the lumped-mass
finite-volume analogue of the finite-element representation in
\citet{lindgren2011spde}. Because it uses the physical cell geometry, it accounts
for the shrinking east--west spacing towards the poles. It can also be restricted to
land cells by deleting rows and columns, and it wraps around in longitude for global
analyses. The continuum relations between $(\kappa, \tau)$ and $(\rho, \sigma)$
hold on the lattice only when $\rho$ spans several cells. A fitted range below two
or three cells (100--150\,km) indicates a field that the $0.5^\circ$ lattice does
not resolve; its $\hat\sigma$ is then not the marginal standard deviation, and we
interpret such fields as cell-level effects.

\begin{remark}\label{rem:boundary}
The discretisation imposes reflecting (Neumann) conditions on the boundary of the
lattice, which roughly doubles the marginal variance of boundary cells
\citep[Section~\ref{sec:sim-engine}; see][for the effect of boundary conditions on
SPDE fields]{khristenko2019boundary}. For a land grid the boundaries are the
coastlines. Extending the lattice into the sea removes the inflation but lets
correlation leak across narrow seas, for example between Yemen and Somalia across the
Gulf of Aden, which is substantively wrong for conflict risk. We therefore use the
barrier model of \citet{bakka2019barrier}: the SPDE is written as
\[
u(s) - \nabla\cdot\Bigl\{\tfrac{r(s)^2}{8}\nabla u(s)\Bigr\} = r(s)\sqrt{\pi/2}\,\sigma\,\mathcal{W}(s),
\]
which reduces to \eqref{eq:spde} with $\kappa = \sqrt{8}/\rho$ when
$r(s) \equiv \rho$. The lattice is extended by six rings of cells (3$^\circ$)
without data. Sea cells, identified from the Natural Earth country polygons, get
range $r(s) = f\rho$. Land cells without data, such as those outside a regional
analysis box, keep the range $\rho$, so that the field is not artificially
decorrelated at the edge of the region. The finite-volume form is
$\bK\bu = \bm{D}\bm{\epsilon}$ with $\bK = \bC + (\rho^2/8)\bG_m$, where $\bG_m$
has edge weights $2m_im_j/(m_i + m_j)\cdot\ell_{ij}/d_{ij}$, $m_i = r_i^2/\rho^2$
(the harmonic mean, the usual two-point flux choice for a discontinuous
coefficient), and $\bm{D} = \mathrm{diag}\{r_i\sigma(\pi a_i/2)^{1/2}\}$, so that
$\bQ_s = \bK\bm{D}^{-2}\bK$. Table~S2 of the Supplementary Material reports the marginal variance
at the coast and the correlation across straits of two to six cells for several
choices of $f$, at $\rho = 400$\,km. At this range, $f = 0.1$, the value of
\citet{bakka2019barrier}, gives a barrier range of 40\,km, less than one cell, and
the coastline behaves almost like a reflecting boundary. We use $f = 0.5$. At
$\rho = 400$\,km it keeps the coastal variance within about 30\% of nominal and the
correlation across straits of four cells or more (about 220\,km) below 0.04, while
a strait of two cells still has a correlation of 0.16. Straits narrower than one
cell are not resolved at this grid resolution. These figures depend on $\rho$,
because the barrier range $f\rho$ and the buffer of six rings scale differently with
it. For the fitted ranges of $w$ in the application, which are well below 400\,km,
the barrier is stronger than in Table~S2 of the Supplementary Material. For ranges of 800\,km or
more, which the SVC fields of Section~\ref{sec:methods-svc} can reach, $f\rho$
exceeds the buffer, so the coastal variance inflation returns and correlation leaks
across the wider straits. The stored real-data fits use the sea cells of the country
polygons as barriers.
\end{remark}

\paragraph{Priors.} We use penalised-complexity priors
\citep{simpson2017pc, fuglstad2019pc}: $P(\rho < \rho_0) = 0.05$ and
$P(\sigma > \sigma_0) = 0.01$, with $\rho_0 = 100$\,km (about two grid cells) and
$\sigma_0 = 1$. The PC prior of \citet{fuglstad2019pc} is derived for Mat\'ern
fields on $\mathbb{R}^2$; we use it as an approximation on the sphere, which is
adequate for ranges far below $R$. The country and cell effects have PC priors with
$P(\sigma_u > 1) = 0.01$ and $P(\sigma_e > 1) = 0.01$.

\subsection{Space-time persistence}\label{sec:methods-st}

The space-time term is a separable field that evolves as a first-order
autoregression of Mat\'ern innovations,
\[
\bv_\tau = \phi\, \bv_{\tau-1} + \sqrt{1-\phi^2}\,\bm{\epsilon}_\tau,
\qquad \bm{\epsilon}_\tau \sim \Normal\{\bm{0}, \bQ_s(\rho_v, \sigma_v)^{-1}\},
\]
\citep{cameletti2013pm10}, so that $\bv$ has precision
$\bQ_{\mathrm{AR}}(\phi) \otimes \bQ_s(\rho_v,\sigma_v)$, where
$\bQ_{\mathrm{AR}}$ is the tridiagonal precision of a unit-variance AR(1) process.
To control cost, $\tau$ indexes blocks of $\Delta$ weeks, and
$v(s,t) = v_{s,\lceil t/\Delta \rceil}$. We use $\Delta = 1$ and $\Delta = 4$ in the
simulations and $\Delta = 8$ in the application. The field $v$ is defined on the
observed land cells only, so it cannot link cells across the sea, and it has
reflecting boundaries at the edges of the observed region. In the models with $v$,
the static field $w$ uses the same lattice as elsewhere: the observed cells, a
buffer of six rings and sea barriers with $f = 0.5$ (Remark~\ref{rem:boundary}). The
prior on
$\operatorname{atanh}\phi$ is $\Normal(0, 1.5^2)$. Non-separable alternatives exist
\citep{lindgren2024diffusion} but are not needed for the questions asked here.

\subsection{State dependence versus persistent heterogeneity}\label{sec:methods-persistence}

Recent event counts dominated the importance rankings of the original boosting
analysis \citep{menicali2026conflicts}. Part of that dominance was an artefact of
duplicated rows (Section~\ref{sec:app-xgb}), but after correction recent history
remains the strongest single predictor. This has several explanations, which
model~\eqref{eq:model} separates only partially:
\begin{enumerate}
\item \emph{state dependence}: recorded events raise the probability of further
recorded events in the same place, through retaliation, mobilisation or the arrival
of armed groups;
\item \emph{persistent heterogeneity}: some places are simply riskier ($u$, $w$, $e$ or $v$
is large there), and past events are informative only as a proxy for that latent
risk \citep{heckman1981heterogeneity};
\item \emph{reporting dynamics}: media attention and reporting networks themselves
persist, so a place that has been reported on is more likely to be reported on
again, whatever happens there \citep{weidmann2016closer}.
\end{enumerate}
All three produce strong autocorrelation in $Y_{st}$ and a large predictive
contribution of $h_{st}$. In the conflict literature the distinction appears as
the question of whether apparent contagion between neighbouring places reflects
spillover or shared conditions \citep{buhaug2008contagion}. Self-exciting point
process models address it by decomposing intensity into a background and a
triggering component
\citep{mohler2011self, mohler2013multisource, porter2012hurdle, clark2018selfexciting,
reinhart2018review, jun2024hawkes}.

In model~\eqref{eq:model}, $\delta$ is the effect of recent recorded history
conditional on the latent persistence $u + w + e + v$, and it is identified only
under the assumed latent structure. It is conditional on heterogeneity at the scales
that are modelled. Persistent heterogeneity that none of the latent terms can
represent loads onto $\delta h_{st}$ and inflates $\hat\delta$. This is why the
cell effect $e$ is included: without it, time-invariant risk that varies between
adjacent cells could only be represented by $h_{st}$.
The space-time field $v$ is resolved at steps of $\Delta$ weeks, and $h_{st}$ covers the previous four weeks. When the two time scales are
similar, separating $\delta$ from $v$ relies on the AR(1) form, so we report
$\delta$ for several $\Delta$ (Section~\ref{sec:sim-persistence}). We therefore read
$\hat\delta$ as conditional dependence of recorded events, not as a causal contagion
effect. To limit the initial-conditions problem of dynamic random-effects models
\citep{wooldridge2005initial}, we drop weeks whose four-week history is incomplete
(the first weeks of each cell and of each country's ACLED coverage). We also add each
cell's history at its first usable week as a covariate, following
\citet{wooldridge2005initial}. This is a limited version of that device. A single
four-week window is a noisy proxy for a cell's initial state, and its coefficient
does not vary in time, whereas under the AR(1) form of $v$ the dependence of later
weeks on the initial state decays with time. The device therefore only partly
removes the initial-conditions bias.

\subsection{Spatially varying climate effects}\label{sec:methods-svc}

To assess whether climate effects differ across biomes and how much variation
remains within them, the effect of anomaly $k$ is decomposed as
\begin{equation}\label{eq:svc}
\beta_k(s) = \gamma_{k,b(s)} + \eta_k(s),
\qquad \bmeta_k \sim \Normal\{\bm{0}, \bQ_s(\rho_k, \sigma_k)^{-1}\},
\end{equation}
a spatially varying coefficient model \citep{gelfand2003svc,dambon2021svc} whose
mean is pooled by biome. The biome means $\gamma_{k,b}$ are fixed effects with the
vague prior of Section~\ref{sec:methods-model}. The field $\bmeta_k$ has PC priors
with $P(\rho_k < 100\,\mathrm{km}) = 0.05$ and $P(\sigma_k > 0.5) = 0.01$, and
$\sigma_k$ measures heterogeneity within biomes. The design of $\bmeta_k$ is the
sparse matrix with entry $z_{k,st}$ in row $(s,t)$ and column $s$.

Three summaries of a biome-specific effect should be distinguished:
\begin{itemize}
\item \emph{Fixed biome slopes.} A model with fixed biome-specific slopes and no
$\bmeta_k$ estimates an information-weighted average,
$\sum_{s \in \mathcal{S}_b}\omega_s\beta_k(s) / \sum_{s \in \mathcal{S}_b}\omega_s$,
with $\omega_s \approx \sum_t z_{k,st}^2\pi_{st}(1-\pi_{st})$. This average puts
most weight on the high-risk cells of the biome.
\item \emph{Biome means.} Model~\eqref{eq:svc} targets the biome mean $\gamma_{k,b}$,
the mean of the process that generated the biome's coefficients. This is a
superpopulation quantity. Because $\bmeta_k$ is not constrained to average zero
within biomes and biomes are large contiguous regions, $\gamma_{k,b}$ is separated
from the biome average of $\bmeta_k$ mainly when $\rho_k$ is small relative to the
size of the biome. When $\rho_k$ is comparable to biome size, the separation rests
partly on the prior. A nearly constant component of $\bmeta_k$ is then confounded
with a common shift of all $\gamma_{k,b}$. The intervals for $\gamma_{k,b}$ are
conditional on $(\hat\rho_k, \hat\sigma_k)$ (Section~\ref{sec:methods-inference}),
and in this situation the uncertainty in these hyperparameters may be the largest
source of error. They should be read with this in mind.
\item \emph{Median SHAP value.} The summary used in the SHAP-based analysis of
\citet{menicali2026conflicts}, the median SHAP value within a biome, differs from
both. The following identities are on the logit scale of the model. They hold for
interventional SHAP values when the features are independent. The path-dependent
TreeSHAP of XGBoost \citep{lundberg2020trees}, which we use, satisfies them only
approximately, and features correlated with biome, such as distance from the
equator, can absorb part of an interaction. For a feature with
an additive effect $f_k(z)$, the SHAP value is $f_k(z) - \E f_k(Z_k)$, so the median
within a biome is $\operatorname{median}_{(s,t):\, b(s) = b}\{f_k(z_{k,st}) - \E
f_k(Z_k)\}$. For a linear effect $\beta_k z$ this is $\beta_k$ times the median of
$z_{k,st} - \E z_k$ in the biome. It reflects where the biome's covariate values lie
relative to the global mean as much as the strength of the effect. When the effect
is a biome-by-feature interaction, SHAP also splits the interaction between the
feature and biome \citep{kumar2020problems,aas2021explaining}. A pure interaction
$g(b)z$, with $B$ independent of $Z$, contributes
$\{g(b) + \E g(B)\}(z - \E z)/2$ to the SHAP value of $z$. Within biome $b$, the
least-squares slope of the SHAP value of $z$ on $z$ therefore targets
$\{g(b) + \E g(B)\}/2$, which is pulled towards the average effect across biomes. It
equals $g(b)/2$ only when $\E g(B) = 0$. The SHAP value with the interaction credited
fully to $z$, $\Phi_{zz} + 2\sum_{j \neq z}\Phi_{zj}$ with $\Phi$ the SHAP
interaction values \citep{lundberg2020trees}, equals $g(b)(z - \E z)$ in this case,
so its slope within biome $b$ targets $g(b)$ whatever the mean of $g$.
\end{itemize}

\subsection{National indicators and spatial confounding}\label{sec:methods-confounding}

National indicators $\bz_{ct}$, such as GDP per capita, the EPI or the ecological
footprint, take one value per country and year. GDP per capita is computed at the
national level (Section~\ref{sec:methods-model}), so that it is constant within each
country and year like the other indicators. We decompose each into a
between-country component, its country mean over the training years
$\bar z_c$, and a within-country component $z_{ct} - \bar z_c$, with separate
coefficients (a Mundlak decomposition). Year effects absorb trends common to all
countries. The between coefficients are identified from contrasts between the
$n_c$ countries only, so their effective sample size is of order $n_c$, not the
number of cell-weeks. Modelling cell-weeks as conditionally independent, or
bootstrapping spatial blocks that cut across countries, therefore produces intervals
that are far too narrow. The within coefficients rest on the few years of change
within each country. A country with no training-period value of an indicator gets
the mean between component and a value of one for that indicator's missingness
indicator (Section~\ref{sec:methods-model}). The missingness indicator is itself
country-level, so its coefficient also rests on few contrasts.

If the indicators are associated with unobserved country-level risk, so that
$\E\{u_c + \bar w_c \mid \bar z_c\} \neq 0$ with $\bar w_c$ the country average of
$w$, then the between coefficients are confounded. An indicator of environmental
performance that tracks state capacity, or reporting capacity, is the leading
example. This is the spatial confounding problem
\citep{reich2006effects, hodges2010adding, paciorek2010importance}. Adding a spatial
random effect changes the estimate of a spatially smooth covariate's coefficient,
and in general neither the estimate with the random effect nor the one without it is
unbiased for a causal effect \citep{khan2022restricted, zimmerman2022deconfounding}.
Recent work has made this more precise. Generalised least squares with a spatial
random effect can be consistent when the covariate has non-spatial variation
\citep{gilbert2025consistency}. Confounding in the data-generating process should be
distinguished from the behaviour of the analysis model \citep{khan2026rethinking}.
The bias depends on how the covariate and the confounder are spread across spatial
scales \citep{guan2023spectral, dupont2026demystifying}. A country-constant indicator
has no within-country spatial variation, so these results give little reassurance
here.

We report the following estimators and treat their disagreement as a measure of how
fragile the national-level conclusions are:
\begin{enumerate}
\item model~\eqref{eq:model} without latent effects;
\item model~\eqref{eq:model} with $\bu$ and $\bm{e}$;
\item model~\eqref{eq:model} with $\bu$, $\bw$ and $\bm{e}$;
\item a \emph{spatial+} estimator \citep{dupont2022spatial}: model~\eqref{eq:model}
with $\bw$ and $\bm{e}$ but \emph{without} $\bu$, in which each between component is
replaced by its residual after an area-weighted projection, over the observed land
cells, on the $m$ smoothest eigenvectors $\bm{\psi}$ of
$\bG\bm{\psi} = \lambda\bC\bm{\psi}$.
\end{enumerate}
These eigenvectors are the lattice analogue of the thin-plate eigenvectors of
\citet{dupont2022spatial}. A country-constant indicator varies only by steps at
borders. After the projection, the residual keeps these steps, minus a smooth
function of the eigenvectors. In a model without $\bu$, the smooth variation is left
to $\bw$ and the border steps identify the coefficient. In a model with $\bu$, the
border steps are country-level jumps, which $\bu$ can absorb. The coefficient would
then be identified only through the shrinkage of $\bu$ under its exchangeable prior,
which is independent of the indicator. That is the random-effects exogeneity
assumption we have just called implausible. This is why (iv) omits $\bu$. Without
$\bu$, however, country-level risk that is not spatially smooth remains a possible
confounder, and the projection also gives the residual some within-country variation
that is a function of the smooth eigenvectors. We therefore treat (iv) as a
sensitivity analysis and not as a preferred estimate.

The default $m^*$ is the number of eigenvalues below
$\kappa_w^2 = 8/\hat\rho_w^2$, the eigenvectors that vary on scales longer than the
fitted range of $w$. This number can be large. By Weyl's law, for a region of area
$|D|$ it is approximately $|D|\kappa_w^2/(4\pi) = 2|D|/(\pi\hat\rho_w^2)$. If
$\hat\rho_w$ is close to the ranges of the hybrid fields of
Section~\ref{sec:app-hybrid} (264\,km for conflict and 68\,km for unrest, on 13{,}719
cells), $m^*$ is of the order of a few hundred for conflict and several thousand for
unrest. The eigenvectors near the cut-off then vary on scales comparable to the size
of many countries, and most of the between-country variation of a country-constant
indicator is projected out. We therefore report $m = 10$, $m = 20$ and $m^*$, and
state $m^*$ for each fit \todo{values of $m^*$ from the rerun of fit\_glmm.py, and
whether the eigenvalue computation reached $m^*$}. The restricted spatial regression
of \citet{reich2006effects} is not used, because its intervals are known to be
anticonservative \citep{hanks2015restricted, khan2022restricted}.

\subsection{Hybrid models: boosting offsets with latent effects}\label{sec:methods-hybrid}

To retain the flexibility of gradient boosting for the many weakly informative
covariates, we also fit a semiparametric version of~\eqref{eq:model}. Its fixed part
is replaced by the logit $\hat f(\bx_{st})$ of an XGBoost classifier
\citep{chen2016xgboost}, which enters as an offset:
\begin{equation}\label{eq:hybrid}
\logit \pi_{st} = \hat f^{(-k(s))}(\bx_{st}) + \beta_0 + u_{c(s)} + \omega(s),
\end{equation}
where $\omega$ is either the Mat\'ern field $w$ or independent cell effects. The
offset is \emph{cross-fitted by country}. Countries are assigned at random to $K = 5$
folds. For each fold $k$, $\hat f^{(-k)}$ is trained on the 2018--2022 rows of the
other folds, with its number of trees chosen by early stopping on 2022 after training
on 2018--2021. It then provides the offset for every row of fold $k$, in both the
training and test periods. The offset and the latent effects are fitted on the
Africa and Middle East box only. Trees that receive national indicators can identify
countries, and through population, biome and distance from the equator, approximate
locations. Cross-fitting by country ensures that the offset of a cell was learned
without any data from its own country, so that $u$ does not count that country's
data twice. It does not exclude data from neighbouring countries. The offset of a
cell near a border can be learned from the outcomes of cells across the border, and
$\omega$ is estimated from the residuals of those same cells. Regional hotspots that
cross borders can therefore enter both the offset and $\omega$. The fitted $\omega$
is a residual field. Its range and standard deviation are not comparable with those
of $w$ in model~\eqref{eq:model}. The offset enters with its coefficient fixed at
one, so fold models whose logits are over- or under-dispersed are not rescaled; $u$
absorbs shifts between folds but not differences in scale. Inference for $u$ and
$\omega$ is conditional on the offset, whose estimation uncertainty is not
propagated, so their posterior uncertainty is understated. Cross-fitting is used here
in the loose sense of sample splitting, not with an orthogonal score as in
\citet{chernozhukov2018dml}.

The price of cross-fitting is that the offset is weaker than a boosting model trained
on all locations. The cross-fitted XGBoost model never sees the country it predicts,
while $u$ and $\omega$ are estimated from that country's training data. A gain of the
hybrid model over the cross-fitted offset alone is therefore not evidence that the
latent effects add anything. The hybrid models are compared, on identical test rows,
with the corrected boosting model trained on all locations, with the location-aware
boosting baseline, with each other ($\omega$ as a field or as independent cell
effects) and with the persistence baseline (Section~\ref{sec:methods-diagnostics}).
The offset uses the same rows, features and hyperparameters as the corrected
boosting model (Section~\ref{sec:app-xgb}). It is trained without class weights, so
it is on the log-odds scale.

Model~\eqref{eq:hybrid} differs from GPBoost \citep{sigrist2022gpboost,
sigrist2023lagaboost}, which alternates boosting steps with latent Gaussian model
updates: here the boosting fit is fixed before the latent model is fitted. This
makes the model cheap enough for the global data and allows the same Laplace
machinery to be used. Random forests with spatially correlated errors
\citep{saha2023rfgls} are a further alternative.

The boosting analysis of \citet{menicali2026conflicts} trained XGBoost with positives
weighted by $\omega_+ = (n_0/n_1)^{1/2}$. For the unrestricted minimiser of the
logistic loss, this shifts the fitted log-odds by $\log\omega_+$
\citep{king2001rare}. We therefore prior-correct the predictions of the weighted
models, $\hat f \leftarrow \hat f - \log\omega_+$, before probabilistic scoring. For
boosted trees the correction is only approximate, because the weights also act on
the leaf penalty, on the minimum leaf weight, which is defined on weighted Hessians,
and on early stopping. The scores of the weighted models therefore include any
remaining miscalibration. Class reweighting does not improve discrimination and
miscalibrates probabilities \citep{vandengoorbergh2022harm}, and the corrected
boosting model is trained without it (Section~\ref{sec:app-xgb}).

\subsection{Inference}\label{sec:methods-inference}

For fixed hyperparameters $\btheta$, the posterior mode
$\bm{\xi}^* = \arg\max_{\bm{\xi}}\{\log p(\bm{y}\mid\bm{\xi}) - \tfrac12 \bm{\xi}^\top \bQ \bm{\xi}\}$
is found by damped Newton iterations,
\[
\bm{\xi} \leftarrow \bm{\xi} + \lambda\,(\bQ + \bA^\top \bm{W} \bA)^{-1}
\{\bA^\top(\bm{y} - \bm{\pi}) - \bQ\bm{\xi}\},
\qquad \bm{W} = \mathrm{diag}\{\pi_{st}(1-\pi_{st})\},
\]
with step halving. The Laplace approximation to the marginal likelihood is
\begin{equation}\label{eq:laplace}
\log \tilde p(\bm{y}\mid\btheta) = \log p(\bm{y}\mid\bm{\xi}^*)
- \tfrac12 \bm{\xi}^{*\top}\bQ\bm{\xi}^* + \tfrac12\log|\bQ| - \tfrac12\log|\bQ + \bH|,
\qquad \bH = \bA^\top\bm{W}^*\bA .
\end{equation}
The hyperparameters are set to
$\hat\btheta = \arg\max\{\log\tilde p(\bm{y}\mid\btheta) + \log p(\btheta)\}$ by
L-BFGS-B with finite-difference gradients on the internal scale (logarithms of ranges
and standard deviations, $\operatorname{atanh}\phi$). Each evaluation warm-starts
Newton at the previous mode. Given $\hat\btheta$, the latent posterior is approximated
by $\Normal\{\bm{\xi}^*, (\bQ + \bH)^{-1}\}$. Samples reuse the factorisation of
$\bQ + \bH$ already computed at the mode, through the perturbation identity
\[
\bm{\xi} = \bm{\xi}^* + (\bQ + \bH)^{-1}\bigl(\bm{R}^\top\bm{\epsilon}_1 + \bA^\top\bm{W}^{*1/2}\bm{\epsilon}_2\bigr),
\]
where $\bm{R}^\top\bm{R} = \bQ$ is block-wise explicit, for instance
$\bm{R} = \bm{D}^{-1}\bK$ for a Mat\'ern block and a Kronecker product of a
bidiagonal AR(1) factor with that matrix for $\bv$. This is an empirical Bayes
version of INLA \citep{rue2009inla} and is similar in spirit to TMB
\citep{kristensen2016tmb}. Laplace approximations are generally accurate for
binary GLMMs but can understate uncertainty when there are few observations per
latent variable and events are rare \citep{fong2010glmm, ferkingstad2015improving}.
The intervals also ignore the uncertainty in $\hat\btheta$.
Section~\ref{sec:sim-engine} therefore checks coverage at the event prevalence of
the data, and Section~\ref{sec:sim-persistence} checks it for the space-time model,
where each latent value is informed by only $\Delta$ Bernoulli observations.

\paragraph{Computation.} All matrices are sparse. Every Newton step and every
log-determinant requires a sparse factorisation of $\bQ + \bH$, whose sparsity
pattern does not depend on $\btheta$. We use a sparse LU factorisation in symmetric
mode (SuperLU). For matrices of dimension above 2{,}000, the fill-reducing ordering
is a nested dissection computed by METIS \citep{karypis1998metis}; it is computed
once and reused, but the factorisation itself is recomputed at every Newton step and
every evaluation of~\eqref{eq:laplace}. If METIS is not available, the code falls
back to multiple minimum degree on $\bm{A}^\top + \bm{A}$. For the space-time model
the nested-dissection ordering reduced factorisation time by a factor of 7 to 14
compared with that fallback. An LU factorisation stores both triangular factors, so
it needs about twice the memory of a Cholesky factor of the same matrix. Newton
iterations stop after at most 50 steps, and a value of $\btheta$ at which the
factorisation fails is given a very low objective value. Both events can make the
optimiser stop at a poor point, so they should be recorded for each fit
\todo{report Newton non-convergence and failed evaluations for the real-data
fits}. Since $\bH = \bA^\top\bm{W}\bA$ aggregates over weeks, the
cost of the spatial and SVC blocks is driven by the number of cells, not by the
number of cell-weeks. A global $0.5^\circ$ land grid gives about $6\times10^4$
latent values per spatial field. The space-time block has $nT_v$ latent values,
which is why it is fitted on a regional subset. The fixed-effect and country columns
add dense rows to $\bQ + \bH$; with tens of fixed effects their cost is small.

\paragraph{Prediction.} For a new cell-week with design row $\bm{a}$ and
offset $o$, the posterior predictive probability is
$\E\{\expit(o + \bm{a}^\top\bm{\xi})\} \approx \expit\{m/\sqrt{1 + \pi s^2/8}\}$, where
$m$ and $s^2$ are the posterior mean and variance of $o + \bm{a}^\top\bm{\xi}$; $s^2$
is estimated from posterior samples. Test years beyond the training period take the
year effect of the last training year.

\subsection{Diagnostics and evaluation}\label{sec:methods-diagnostics}

\paragraph{Residual diagnostics.} Binary residuals are uninformative one at a
time, so we aggregate over the test weeks within cells. The standardised cell
residual is
\[
r_s = \sum_{t} (y_{st} - \hat\pi_{st}) \Big/ \Bigl\{\sum_t \hat\pi_{st}(1-\hat\pi_{st})\Bigr\}^{1/2},
\]
with the sums over the test period 2023--2024. All residual diagnostics reported in
the paper are computed in this way for every model, including the cross-fitted
offset, on the same test rows. Residuals of in-sample latent-model fits are shrunk
and spatially dependent through the smoother, and we do not use training-period
residuals as evidence.

Even for a correct model, the test residuals are neither standard normal nor
independent across cells. The predictions $\hat\pi_{st}$ carry the estimation error
of the training period, for example in $\hat w(s)$. A logit error $\epsilon_s$
contributes about $\epsilon_s\{\sum_t \pi_{st}(1-\pi_{st})\}^{1/2}$ to $r_s$, which
grows with the number of test weeks, and the errors of a smoothed field are
positively correlated between neighbours. The diagnostics below are therefore used
to compare models on the same rows, not as tests against a null of independence. We
report:
\begin{itemize}
\item Moran's $I$ \citep{moran1950notes} of $\{r_s\}$ for rook adjacency. We also give
a permutation $p$-value from $B = 999$ permutations, which is bounded below by
$1/(B + 1)$. Its null hypothesis of exchangeable residuals does not hold even for a
correct model, so it is descriptive only;
\item the empirical semivariogram of $\{r_s\}$ against great-circle distance;
\item the correlation of $r_s$ between neighbouring cells in the same country and
across a border.
\end{itemize}
A larger within-country correlation can point to effects at the country level, of
occurrence or of reporting. For the cross-fitted offset it also has a mechanical
source. Neighbours in the same country are always predicted by the same fold model,
while neighbours across a border are usually predicted by different fold models,
whose differences create a country-level residual component. For this model we
therefore also split the across-border pairs into pairs whose two countries are in
the same fold and pairs in different folds.

\paragraph{Out-of-sample evaluation.} Models are trained on 2018--2022 and
evaluated on 2023--2024, matching the original analysis. All models are compared on
identical test rows. We report the logarithmic score and the Brier score, both
strictly proper \citep{gneiting2007strictly}, and the area under the
precision--recall curve, computed as average precision \citep{boyd2013area}. The
latter is more informative than the ROC curve at the low prevalence of the data
\citep{saito2015precision}. Paired differences in log score and Brier score are
reported with two standard errors: one clustered by cell and one clustered by
country, with the number of clusters. The cell-clustered errors allow for dependence
over time within a cell but not between cells. The latent fields and the national
indicators induce dependence at regional and national scales
(Section~\ref{sec:methods-confounding}), so the country-clustered errors are the
more reliable of the two, although they rest on few clusters. Since prediction in
the same places is dominated by persistence \citep{mueller2022hard}, every model is
compared with a persistence baseline that ranks cell-weeks by the number of events in
the previous four weeks. We also report all scores on \emph{onset} rows: test
cell-weeks in cells that had no event of that outcome in the previous 52 weeks. On
these rows recent history carries no information, so they show how well a model
anticipates events in places that have been quiet.

We also compare every model with a \emph{location-aware} boosting baseline that adds
each cell's training-period event rate as a feature. For the fit used to choose the
number of trees by early stopping on 2022, this rate is computed from 2018--2021
only, so it uses no 2022 outcomes. For the final fit on 2018--2022 it is computed
leaving out the row's own year, and for the test rows it is the rate over
2018--2022. This baseline is the boosting counterpart of a cell intercept, and the
latent fields have to improve on it to be worth their cost.
